\documentclass[11pt]{article}
\usepackage[margin=1in]{geometry}
\usepackage{amsmath,amssymb}
\usepackage{booktabs}
\usepackage{array}
\usepackage{hyperref}
\hypersetup{colorlinks=true,linkcolor=black,citecolor=black,urlcolor=blue}

\title{\textbf{Dataset Collision-Checking Workflow}\\[2pt]\large Detection and Attribution of AI-Generated and Human Written
Phishing Emails from Legitimate Emails}
\author{Mannan Bhardwaj -- 25BYB0013 \\ Bhoumik Singh Rajput -- 25BYB0120 \\[4pt]
School of Computer Science Engineering and Information Systems}
\date{}

\begin{document}
\maketitle

\section{Purpose}
Before any Naive Bayes sub-model is trained on \texttt{training\_set.csv}, the dataset must be checked
for \emph{collisions}: cases where the row-level bookkeeping (identifiers, class labels, and the
train/holdout split) is internally inconsistent in a way that would silently corrupt training or inflate
reported accuracy. This is distinct from the earlier statistical validity checks (Kruskal-Wallis,
Chi-Square, the TF-IDF baseline classifier), which ask \emph{whether the features carry real signal}.
The workflow below instead asks \emph{whether the dataset's own bookkeeping is self-consistent} --- a
necessary precondition that has to hold regardless of which features are eventually used. It was run
against the corrected (recipient-identity-leak-free) versions of \texttt{training\_set.csv} and
\texttt{holdout\_set.csv}.

\section{The four collision classes}
Each email in the dataset is assigned an identifier \texttt{id}, computed as a truncated (40-bit) MD5
hash of its body text, prefixed by its source corpus. Four distinct failure modes are checked, in order
from most mechanical to most subtle.

\subsection*{1. ID collision (hash collision)}
Two rows share the same \texttt{id} value but have \emph{different} body text. Since \texttt{id} is a
hash of the body, this would only happen from a genuine MD5 collision or from a pipeline bug (e.g.\ the
hash computed before a text-normalizing step was applied inconsistently). At the dataset's row count
($\sim$13{,}000 total), the expected number of genuine hash collisions under a 40-bit hash is
approximately
\[
\frac{n^2}{2 \cdot 2^{40}} \approx \frac{13000^2}{2 \cdot 2^{40}} \approx 4 \times 10^{-5},
\]
i.e.\ essentially zero. A nonzero count here should be treated as a pipeline bug, not bad luck.

\subsection*{2. Cross-class label collision}
The exact same body text appears more than once with \emph{different} class labels (e.g.\ the identical
email tagged both \textit{Legitimate} and \textit{Human-Phishing} in different rows). This is the most
serious possible failure: it is a direct contradiction in the ground truth, and any model trained across
such rows is being given mutually exclusive answers for the same input.

\subsection*{3. Train/holdout overlap}
Any row --- checked both by \texttt{id} and, independently, by raw body text --- appears in \emph{both}
\texttt{training\_set.csv} and \texttt{holdout\_set.csv}. This would let information used for training
leak into the holdout set, inflating any accuracy reported against it.

\subsection*{4. Duplicate content (within a file)}
The same body text appears more than once within a single file. Unlike the previous three, this is not
automatically a failure --- a phishing template genuinely sent to many recipients produces legitimately
repeated content --- but it is tracked as a regression guard against unintended re-introduction of exact
duplicates (the kind removed during corpus cleaning).

\subsection*{5. Near-duplicate template reuse (informational)}
Exact-match checks (1--4) miss emails that are identical \emph{except} for a date, dollar amount, or
name. These are detected by normalizing each body (lower-casing, collapsing all digits to a single
placeholder character, stripping punctuation, collapsing whitespace) and then checking for repeats under
that normalized form. This is not treated as a hard failure --- template reuse is a realistic property of
a phishing corpus --- but any near-duplicate cluster that spans \emph{more than one class} is flagged
specifically, since that would be equivalent in severity to a cross-class label collision.

\section{Results on the current dataset}
Table~\ref{tab:results} reports the outcome of running this workflow against the corrected
\texttt{training\_set.csv} (4{,}491 rows) and \texttt{holdout\_set.csv} (8{,}465 rows).

\renewcommand{\arraystretch}{1.3}
\begin{table}[h]
\centering
\begin{tabular}{@{}p{7.2cm}rr@{}}
\toprule
\textbf{Check} & \textbf{Training set} & \textbf{Holdout set} \\
\midrule
ID collisions (different content, same id) & 0 & 0 \\
Cross-class label collisions & 0 & 0 \\
Duplicate-content groups (exact match) & 0 & 0 \\
Near-duplicate groups (normalized match) & 45 (95 rows) & 38 (133 rows) \\
Near-duplicate clusters crossing class boundary & 0 & 0 \\
\bottomrule
\end{tabular}
\caption{Collision-check results. The four hard-failure checks (rows 1--3 and the cross-file row below)
are all 0.}
\label{tab:results}
\end{table}

\vspace{-6pt}
\begin{table}[h]
\centering
\begin{tabular}{@{}p{7.2cm}r@{}}
\toprule
\textbf{Cross-file check} & \textbf{Result} \\
\midrule
Train/holdout overlap by \texttt{id} & 0 \\
Train/holdout overlap by body text & 0 \\
\bottomrule
\end{tabular}
\end{table}

All four hard-failure checks pass. The 45 (training) and 38 (holdout) near-duplicate clusters are
template-reuse within a single class only --- manual inspection confirmed 91 of the 95 affected training
rows are Human-Phishing (consistent with template phishing campaigns sent to multiple recipients) and the
remaining 4 are Legitimate; none span class boundaries, so none represent a labeling problem.

\section{Re-running the workflow}
This check should be re-run any time the dataset-building pipeline (\texttt{clean\_and\_merge.py}) is
modified, or before any claim of dataset validity is made in the written paper. It is intentionally kept
separate from the statistical-validity checks (Section~3 of the proposal) because it tests a different
property: internal bookkeeping consistency, not feature informativeness. A pipeline change that alters
row counts, deduplication, or redaction logic should always be followed by re-running both this workflow
and the statistical baseline check before new files are treated as final.

\end{document}
